
We tested whether providing information from randomized trials about
flu vaccine safety influences the beliefs and vaccination decisions of
vaccine hesitant Americans. 
Providing such information is a potentially  welfare-enhancing intervention because 
flu vaccine hesitant individuals hold
pessimistic beliefs about vaccine side effects, and these beliefs appear 
important for their vaccination decisions.  However, flu vaccine hesitancy
is associated with mistrust of government's and doctor's advice, making 
it unclear whether information can shift beliefs or behaviors.

Our main result is that providing side effect information shift
beliefs and intentions, but belief updating is limited and moderated by
relatively low trust in trials. Providing information from academic and industry 
sponsored trials yielded similar belief updating. Emphasizing representativeness
of study populations for vaccine hesitant individuals did not lead to updating. 
We found mixed evidence that the side effect information increased vaccinations. 
Information from the academic sponsored trial increased intentions to vaccinate,
but information from the industry sponsored trial did not. We lacked the power
to detect effects on actual vaccination. 

While our results suggest side effect information can be an important part of messaging 
to increase vaccination, such messaging is likely best paired with a trustworthy messenger.
Overall trustworthiness of the described studies was not especially high, but effects
were large for participants who found the study trustworthy. 
Prior work such as \citet{larsen2023counter} and \citet{alsan2024experimental} has found
that  a credible messenger increases the persuasiveness of a pro-vaccine message, so
an especially promising direction for future work is to pair an appropriate messenger
with a message of vaccine safety. 